\subsection{Design of Protocols}\label{sec:security_of_cryptographic_protocols:design}

The design of cryptographic protocols comprises several steps, which we illustrate below.

\paragraph*{1. Scenario and Security Properties.} First, the designers of a cryptographic protocol gather detailed information about the scenario and elicit desired security properties by engaging stakeholders. Notably, this step has more to do with \textbf{requirements engineering} rather than with technical aspects --- and is always more complex than it seems. Relevant information and considerations resulting from this step are the list of interacting parties, assets to protect, use cases, regulatory constraints, and the operational environment (e.g., performance constraints on computational resources, latency, and bandwidth).

\warningBlock{Interdependencies among security properties}{Beware that some security properties may have interplays (e.g., integrity vs. sender authenticity vs. ciphertext authenticity) and even conflict with each other (e.g., non-repudiation vs. deniability).}

\paragraph*{2. Assumptions and Threat Model.} Then, the designers explicitly state all assumptions related to trust (e.g., which parties are trusted) and security (e.g., on cryptographic material available to those parties and on cryptographic keys binding --- see \Cref{sec:binding_keys_to_parties}). At the same time, the designers define a threat model (\Cref{sec:cryptographic_security_overview:threat_model}) for enumerating and prioritizing possible attacks (e.g., see \Cref{sec:cryptographic_security_overview:attacks,sec:security_of_cryptographic_protocols:design:attacks}) to the cryptographic protocol present in the considered scenario, thereby ensuring that the protocol addresses realistic risks rather than hypothetical or impractical ones. A good threat model should be realistic, comprehensive, and consistent, and it should be updated regularly to reflect the changing environment and threats.

\exampleBlock{Threat models for cryptographic protocols}{Communication channels are often assumed to be controlled by an attacker. Starting from this assumption, two threat models are commonly used for cryptographic protocols. The \textbf{Dolev-Yao} threat model \cite{dolev_security_1983} is widely used for discussing (and formally verifying) the security of protocols \cite{wazid_authentication_2019}. The Dolev-Yao threat model defines an attacker --- i.e., the Dolev-Yao attacker --- who has full control over communication channels, i.e., the attacker can intercept, modify, inject, delete, and replay any message. Moreover, the Dolev-Yao attacker can execute and interleave multiple executions (or \textbf{runs}) of protocols. However, the Dolev-Yao attacker cannot affect single components, break cryptography (e.g., through cryptanalysis, brute-forcing, or side-channel attacks --- see \Cref{sec:security_of_cryptographic_algorithms:attacks}), or affect the physical world in any way. In other words, the Dolev-Yao threat model assumes perfect cryptography, operating in the symbolic model (\Cref{sec:security_of_cryptographic_protocols:security}). Another widely used threat model in the context of protocols is \textbf{Canetti-Krawczyk} \cite{canetti_universally_2002} which defines an attacker with the same capabilities as the Dolev-Yao attacker; in addition, the Canetti-Krawczyk attacker can also conduct physical attacks to potentially retrieve secret data from cryptographic modules, operating in the computational model (again, see \Cref{sec:security_of_cryptographic_protocols:security}).}

\paragraph*{3. Informal Flow and Algorithms Selection.} Once the context and the objectives are clear, the designers can draft an informal description of the cryptographic protocol in plain language or simple diagrams (e.g., message sequence charts): this amounts at defining the messages (in terms of syntax, semantics, and temporization) and the cryptographic algorithms to employ (e.g., hash functions, symmetric and asymmetric ciphers). Usually, this is a top-down process where details are refined iteratively to balance security properties with performance constraints (actually, the whole design process is often iterative).

\readBlock{The Alice and Bob notation}{Dating back to 1978, the \textbf{Alice and Bob notation} is a simple, intuitive, and popular graphical language for the description and specification of the interactions between system components (e.g., such as cryptographic protocols). Quinn DuPont's and Alana Cattapan's website ``A History of The World's Most Famous Cryptographic Couple'' --- available at the link \url{https://cryptocouple.com/} --- details the major events in the ``lives'' of Alice and Bob, from their birth in 1978 onward.}

\paragraph*{4. Formal Specification and Security Analysis.} Afterwards, the designers translate the informal design using a formal language or pseudocode to avoid ambiguities. Having a formal specification allows for then applying formal methods (\Cref{sec:security_of_cryptographic_protocols:security}) to prove the security of the cryptographic protocol under the defined threat model.

\paragraph*{5. Configuration, Implementation, and Deployment.} Lastly, the designers can define a suitable configuration (e.g., cipher suite choices, key lengths, timestamp windows, error‑handling semantics), choose suitable cryptographic libraries, and deploy the cryptographic protocol.

\remarkBlock{Further points of attention}{Besides the aforementioned 5 steps, the secure design of cryptographic protocols needs to take into account key management (\Cref{sec:key_management}), proper randomness sources,
%TODO (\Cref{future_sec:randomness}), 
clear deployment and usage guidelines, and update and deprecation mechanisms.}

\subsubsection{Best Practices}\label{sec:security_of_cryptographic_protocols:design:best_practices}

Even years after their deployment, many cryptographic protocols were found to have vulnerabilities and weaknesses; learning from these mistakes, practitioners and researchers defined best practices for their design. Unfortunately, there is not a unique list of best practices formally agreed upon by everyone; nonetheless, below we make an effort to review the main (i.e., most important) design best practices to avoid common pitfalls when designing cryptographic protocols.

\exampleBlock{A famous example of a vulnerable cryptographic protocol}{A replay attack on the symmetric version of the Needham-Schroeder key exchange protocol (comprising 5 messages only) went undetected for around 3 years; even more notably, a \gls{mitm} attack on the public key version of the Needham-Schroeder protocol (comprising 7 messages only\footnote{The original version of the public-key Needham-Schroeder protocol in 1978 \cite{needham_using_1978} assumed a trusted server for key distribution within a 7-step protocol. However, later treatments (notably, Lowe's 1995 attack \cite{lowe_attack_1995}) abstract away the server, focusing on the core 3-message exchange.}) was discovered after around 17 years.\footnote{Needham–Schroeder protocol - Wikipedia (\url{https://en.wikipedia.org/wiki/Needham–Schroeder_protocol})}}

\paragraph*{Map Cryptographic Protections with Security Properties.} Map each cryptographic protection (e.g., using a table) to the security properties it provides --- albeit it may seem obvious --- rather than relying on intuition or colloquial terminology. Subtle protocol flaws may derive from mismatches between cryptographic protections and security properties, e.g., assuming that encryption alone yields integrity, that freshness or ephemeral keys imply authenticity, or that \gls{mac} (\Cref{sec:symmetric_cryptography:mac}) provides non‑repudiation. It is also important to rigorously check that the composition of cryptographic protections does cover every required security property.

\paragraph*{Simplicity and Clarity.} Being arguably (almost) always part of the \gls{tcb} of a system (\Cref{sec:programming_languages_and_software:trusted_computing_base}), design cryptographic protocols as simple and minimal as possible. Every feature, message field, or option must have a clear security purpose. However, note that simplicity and clarity do not mean lacking functionality.

\paragraph*{The Human Factor.} Design cryptographic protocols with human comprehension and usability in mind; this includes both end users and implementers. In fact, even the most secure cryptographic protocol can be compromised by user error, misuse, misconfiguration, or misunderstanding. Concretely, use secure defaults and reduce opportunities for incorrect choices and ensure that any required user actions (e.g., confirming an identity, accepting a key, or verifying a fingerprint) are clear and guided. Provide friendly interfaces, meaningful error messages, clear documentation, and, wherever applicable, training materials or educational cues: if users do not understand what a security warning means, they may ignore it. Similarly, if implementers do not understand the specification of a cryptographic protocol, they may create insecure variants.

\paragraph*{Freshness.} Include freshness to ensure that messages are recent and not reused to thwart reflection and replay attacks. Common methods to include freshness in cryptographic protocols are nonces and timestamps: nonces are non-secret random or sequential numbers (either way, the important thing is to avoid collisions) used only once and protect the parties that generated them,\footnote{A nonce protects the party that generated it against replay attacks. For instance, consider a party $A$ that needs to submit a query to a party $B$: $A$ generates a nonce $n_A$ and sends it along with the query to $B$. Then, $B$ replies by sending the answer to the query and the nonce $n_A$ to $A$. Now, $A$ is sure that the answer is ``fresh'', given that it includes the newly generated nonce $n_A$ (assuming no collisions happened and that no attacker can replace nonces in messages). However, $B$ is not protected against replay attacks --- unless $B$ keeps a list of all nonces ever received and checks the list every time it receives a new query to detect replay attacks; however, this kind of protection against replay attacks is often impractical due to the complexity in handling a continuously growing list. Therefore, it is more common for $B$ to generate a nonce $n_B$ as well (or to keep the list of received nonces short by requiring queries to include timestamps and consequently preserve a nonce in the list only if within the allowed clock skew).} while timestamps ensure message timeliness (with an allowable clock skew) for all parties involved. Intuitively, timestamps are tricky to use as they inherently have a time window tolerance (e.g., a few seconds) that requires synchronization among parties and entails a short-lived vulnerability to replay attacks.

\paragraph*{Channel Binding.} Include context information (e.g., identities of sender and recipient, session identifiers, purpose) in each message --- typically, in \gls{aead} authentication tags (\Cref{sec:symmetric_cryptography:ae:aead}) --- to explicit under what circumstances a message is valid (and useless elsewhere) and thwart impersonation, rebinding, reflection, message insertion, and reordering attacks. In other words, implicit context (e.g., assuming the recipient ``knows'' what a message means based on order) should be avoided. Similarly, it is a best practice to ``chain'' cryptographic computations through the injection of context information (e.g., encrypt data with committing \gls{aead} using a symmetric key established from a combination of asymmetric keys and parties' identifiers).

\exampleBlock{Channel binding in real-world protocols}{Modern cryptographic protocols such as \gls{tls} v1.3 and Signal 
%TODO (\Cref{future_sec:popular_protocols:tls,future_sec:popular_protocols:signal})
inject context information in key establishment or in \gls{aead} computations. In fact, channel binding applies also to layered authentication, i.e., whenever a cryptographic protocol runs on top of another (e.g., Signal over \gls{tls}) in a secure channel \cite{rfc5056}.}

\paragraph*{\gls{aead}.} If symmetric cryptography is needed, use (committing) \gls{aead} ciphers such as \gls{aes}-\gls{gcm} and ChaCha20-Poly1305 (\Cref{sec:symmetric_cryptography:ciphers:aes,sec:symmetric_cryptography:ciphers:chacha20}) to thwart ciphertext forgery and tampering attacks.

\paragraph*{\glspl{chf}.} If digests are needed, use robust \glspl{chf} such as SHA-3 and BLAKE2 to thwart collision, pre-image, and length extension attacks (\Cref{sec:hash_functions:cryptographic_hash_function}). Moreover, include context information or a prefix as input to \glspl{chf} to avoid cross-protocol collision issues.

\warningBlock{Using \glspl{chf} for integrity}{If digests are required for integrity, do not ``roll your own'' integrity checks by hashing data and comparing the resulting digests manually --- instead, use \gls{hmac} or digital signatures as appropriate (\Cref{sec:symmetric_cryptography:hmac,sec:digital_signatures}).}

\paragraph*{\gls{pkc}.} If asymmetric cryptography is needed, use robust asymmetric ciphers for encryption (\Cref{sec:asymmetric_cryptography:ciphers}) as the basis for hybrid cryptography (\Cref{sec:hybrid_cryptography}) and robust asymmetric ciphers for digital signatures (\Cref{sec:digital_signatures:ciphers}).

\paragraph*{Ephemeral Keys.} Include ephemeral keys (\Cref{sec:cryptography_basics:keys}) to provide forward secrecy during key establishment (\Cref{sec:key_management:lifecycle:pre}) --- typically when deriving (symmetric) session keys --- and mitigate long-term confidentiality violation attacks.

\exampleBlock{Ephemeral keys in real-world protocols}{The use of ephemeral keys to provide forward secrecy is enforced by default in modern cryptographic protocols such as \gls{tls} v1.3 and Signal. Moreover, Signal allows for providing limited forward secrecy by using short-term asymmetric key pairs.}

\paragraph*{Key Separation or Single Purpose.} Prefer using distinct keys for distinct purposes (e.g., encryption vs. authentication vs. digital signatures vs. different communication directions – client-to-server and server-to-client), even if derived from a common secret through \glspl{kdf} with proper domain separation, e.g., injecting labeled context strings (\Cref{sec:key_management:lifecycle:pre}). In fact, reusing keys across different cryptographic algorithms can weaken their security guarantees against cryptanalysis and amplify the impact of key disclosure \cite{barker_recommendation_2020_p1}. However, note that the same key may sometimes be used to provide multiple security properties (e.g., confidentiality and ciphertext authenticity in \gls{aead}), though almost always in the context of a single cipher.

\remarkBlock{Ciphers not respecting the key separation principle}{There are ciphers which disregard key separation for a more pragmatic design, but only after a very careful trade-off and design analysis --- see \Cref{sec:digital_signatures:ciphers:eddsa} for some examples and more details.}

\paragraph*{Security Over Backward Compatibility.} Do not sacrifice security for backward compatibility --- that is, do not include obsolete cryptographic algorithms or configurations to accommodate outdated systems: prioritizing or even allowing backward compatibility may (re-)introduce vulnerabilities. If legacy devices are present, need to communicate, and cannot be updated, consider offering a separate legacy mode if absolutely necessary (and remember to isolate the legacy mode from the main secure mode).

\exampleBlock{\gls{tls} and backward compatibility}{\gls{tls} v1.2 favored backward compatibility with \gls{tls} v1.1 by considering 45 cipher suites --- inheriting serious vulnerabilities --- while the latest \gls{tls} v1.3 favors security by considering 5 cipher suites only and removing insecure options (such as \gls{rsa} key establishment).}

\paragraph*{Crypto-Agility.} Cryptographic protocols must be capable of swapping out cryptographic algorithms without a complete redesign and with minimal disruption. In fact, cryptographic algorithms deemed robust may become weak due to cryptanalytic breakthroughs or new capabilities (e.g., quantum-computing).
%TODO --- see \Cref{future_sec:post_quantum_cryptography}).
Consequently, include versioned algorithm identifiers (possibly standard) and cipher suites negotiation secured against downgrade attacks \cite{chen_considerations_2025} and be algorithm-neutral in design. 

\remarkBlock{Trade-offs between crypto-agility and complexity}{If necessary (e.g., the cryptographic protocol targets a wide generic audience), crypto-agility would prescribe to offer more options for cipher suites to avoid scenarios in which a single cipher gets broken and leaves no alternative. However, this would add further complexity, increase the \gls{tcb} (\Cref{sec:programming_languages_and_software:trusted_computing_base}), and ultimately enlarge the attack surface. As a unique optimal choice does not exist, the designers of a cryptographic protocol should identify the best trade-off between crypto-agility and complexity on a case-by-case basis.}

\paragraph*{Post-Quantum Cryptography.} If expected by the threat model, design protocols to include --- either during the first design or also in the future --- post-quantum cryptographic algorithms (e.g., by planning for larger keys, message sizes, and digital signatures) to thwart ``harvest now, decrypt later'' attacks.
%TODO (\Cref{future_sec:post_quantum_cryptography}). 

\paragraph*{Standards.} Consider standards (\Cref{sec:standards_and_certifications}) and white papers of reputable organizations (\Cref{sec:appendix:further_resources}) when designing cryptographic protocols to avoid known mistakes and enhance interoperability, security, and unambiguity. Finally, being consistent with standards also helps with compliance.

\paragraph*{Performance Constraints.} Be mindful of performance constraints (e.g., low power, low bandwidth, high latency, limited UI, low battery or CPU) tied to the specific application or environment (e.g., \gls{iot} vs. web vs. mobile vs. automotive) in which a cryptographic protocol will be deployed: cryptographic protocols impractical in their environment will either fail to be adopted or, worse, be implemented in an insecure way to ``make them work''. This may imply prioritizing security properties (e.g., forward secrecy and deniability for instant messaging protocols, authenticity and confidentiality for \gls{iot}) according to the considered threat model. Still, we recommend never compromising core security principles for performance.

\paragraph*{Testing and Verification.} Analyze and test the security of cryptographic protocols by conducting attack simulations or using formal methods and verification tools (\Cref{sec:security_of_cryptographic_protocols:security}). If possible, engage in security audits or invite cryptographers to review the design of cryptographic protocols.

\paragraph*{Ratcheting.} Ratcheting mechanisms
%TODO (\Cref{future_sec:cryptographic_ratchets})
ensure that cryptographic protocols can recover from key compromise over time. Often, an attacker who obtains currently used private or secret keys may decrypt all future messages. In this context, ratcheting mechanisms allow cryptographic protocols to self-heal after compromise --- that is, once the attacker is out of the loop, future messages become secure again.

\paragraph*{Error Handling.} Handle errors due to failures, inconsistent states, and unforeseen situations properly by preventing information leakage and defining clear recovery or abort procedures to ensure predictable and safe behavior. Protocol-level error messages must be uniform and generic, revealing only that an operation failed without exposing details such as whether a decryption failed due to an incorrect padding value, an expired key, or a malformed ciphertext. Error messages in logs may contain more information, but still must not disclose sensitive information (e.g., private or symmetric cryptographic keys).

\subsubsection{Categories of Attacks}\label{sec:security_of_cryptographic_protocols:design:attacks}

In order to design secure cryptographic protocols, one should first have an idea of what the (at least most common) attacks are. Compiling an exhaustive list is difficult --- as each cryptographic protocol may present peculiar vulnerabilities and weaknesses --- but we can identify the following high-level attack categories:

\begin{itemize}
	
	\item \textbf{passive eavesdropping attacks}: the attacker listens to communications without altering them to read unencrypted data or perform traffic analysis (e.g., to identify patterns in timing, frequency, or length of messages);
	
	\item \textbf{\gls{mitm} attacks}: standing between two (or more) communicating parties, the attacker exploits poor authenticity to impersonate one or both parties by intercepting and possibly altering exchanged messages. Intuitively, the legitimate parties are usually not aware of the \gls{mitm} attacker;
	
	\item \textbf{downgrade attacks}: the attacker forces parties engaging in a cryptographic protocol to fall back to outdated or less secure protocol versions and cipher suites, often by interfering during the negotiation or handshake phase\footnote{A handshake in a cryptographic protocol is the initial exchange of messages between two (or more) parties whose goals are usually to authenticate one another, negotiate protocol parameters, and establish fresh secrets.};
	
	\exampleBlock{Real-world downgrade attacks}{Two famous downgrade attacks are the POODLE (Padding Oracle on Downgraded Legacy Encryption)\footnote{SSL 3.0 Protocol Vulnerability and POODLE Attack | CISA (\url{https://www.cisa.gov/news-events/alerts/2014/10/17/ssl-30-protocol-vulnerability-and-poodle-attack})} and the FREAK (Factoring RSA Export Keys)\footnote{FREAK SSL/TLS Vulnerability | CISA (\url{https://www.cisa.gov/news-events/alerts/2015/03/06/freak-ssltls-vulnerability})} attacks on \gls{tls}.}
	
	\item \textbf{replay attacks}: the attacker delays, repeats, or reuses valid protocol messages --- often originated by legitimate parties --- in different and possibly simultaneous runs of cryptographic protocols. For instance, an attacker might capture an authentication token and send it again to impersonate the original sender at a later time. Two sub-categories of replay attacks are \textbf{reflection attacks} --- where the attacker tricks a party into unknowingly responding to its own messages --- and \textbf{rebinding attacks} --- where the attacker reuses messages in different contexts or for different purposes than originally intended;
	
	\item \textbf{oracle attacks}: the attacker takes advantage of normal protocol responses (e.g., error messages, response codes) or their characteristics (e.g., timing differences) to make a legitimate party --- in this context called \textbf{oracle} --- reveal or leak secret data unknowingly, often as a encryption or decryption service. Typically, oracle attacks require multiple interactions;
	
	\exampleBlock{Real-world oracle attacks}{A classic example of oracle attacks are \textbf{padding oracle attacks}: the attacker sends malformed encrypted messages to a server and observes whether the server returns a padding error; by varying the ciphertext and watching the server's behavior, the attacker can gradually deduce the correct plaintext or encryption key. An instance of a padding oracle attack is the Bleichenbacher attack \cite{bleichenbacher_chosen_1998} already discussed in \Cref{sec:cryptographic_security_overview:attack_model}.}
	
	\item \textbf{\gls{dos} attacks}: the attacker aims to violate the availability of a system. In reality, \gls{dos} attacks typically do not concern cryptography, except when a cryptographic protocol is (badly) designed and allows the attacker to repeatedly trigger costly cryptographic computations (e.g., by flooding a server with bogus handshake initiations);
	
	\item \textbf{online guessing attacks}: the attacker conducts guessing attacks by actively interacting with a system, e.g., repeatedly attempting to authenticate using a different password or to decrypt some data by guessing a decryption key. Intuitively, online guessing attacks are relatively easy to detect and mitigate (e.g., with lockout policies and rate limiting --- which however may yield \gls{dos} attacks);
	
	\item \textbf{offline guessing attacks}: the attacker obtains cryptographically-protected data (e.g., an encrypted file, a password hash) and then attempts to break the cryptographic protection offline via a trial-and-error approach. Offline guessing attacks are often regarded as passive attacks and are effective against cryptographic protocols employing low-entropy sources.
	%TODO (\Cref{future_sec:randomness}).
	Brute force, dictionary, and rainbow table attacks are instances of offline guessing attacks (\Cref{sec:cryptographic_security_overview:attacks});
	
	\item \textbf{side-channel attacks}: the attacker relies on leakage of cryptographic computations to infer (portions of) secret data such as seeds, private keys, and internal states. In cryptographic protocols, the most relevant side-channel attacks are remote timing attacks, where the attacker measures how long a party takes to respond (or differences in response times) to infer secret data. Side-channel attacks are thoroughly discussed in \Cref{sec:security_of_cryptographic_algorithms:attacks:sca};
	
	\item \textbf{``harvest now, decrypt later'' attacks}: the attacker records encrypted data hoping to decrypt it in the future when more computational power will likely be available (see the Moore's law in \Cref{sec:cryptographic_security_overview:attacks} and, especially, pos-quantum cryptography).
	%TODO in \Cref{future_sec:post_quantum_cryptography}). 
	Note that in the context of ``harvest now, decrypt later'' attacks, forward secrecy does not help in protecting future data.
	
	\item \textbf{\gls{kci} attacks}: the \gls{kci} attack assumes that an attacker was able to compromise the long-term private key of a party used in key establishment (e.g., such as \gls{dh} --- see \Cref{sec:diffie_hellman}). A natural consequence is that now the attacker can pretend to be that party when communicating with everyone else. However, in cryptographic protocols suffering from \gls{kci} attacks, the attacker can also pretend to be anyone else when communicating with that party (due to the symmetric nature of \gls{dh}). Preventing \gls{kci} attacks generally involves either an interactive authenticity protocol or digitally signing messages (but forcing then non-repudiation, which may not be a desired security property). A cryptographic protocol suffering from \gls{kci} attacks is Signal;
	
	\item \textbf{\gls{uks} attacks}: the attacker manages to trick one party $A$ into believing to have established a shared secret key with another party $B$ while in reality $A$ established the key with a third party $C$. Instead, $B$ correctly believes to have established the key with $A$. Usually, a \gls{uks} attack does not allow the attacker to learn the key; still, the attacker can abuse $A$'s misassociation. \gls{uks} attacks are typically possible when parties' identities are included neither in the \gls{kdf} used to obtain the shared secret key nor in key confirmation messages;
	
	\item \textbf{Truncation attacks}: the attacker drops an arbitrary amount of messages without the communicating parties noticing it.
	
	\exampleBlock{A real-world truncation attack}{A recent example of truncation attack is the \gls{ssh}
	%TODO (\Cref{future_sec:popular_protocols:ssh})
	\textbf{Terrapin attack} \cite{baumer_terrapin_2024}:\footnote{Terrapin Attack (\url{https://terrapin-attack.com/})} by carefully adjusting the sequence numbers during the handshake, the attacker can remove an arbitrary amount of messages sent by the client or server at the beginning of the secure channel, downgrading the connection's security by truncating the extension negotiation message (RFC8308 \cite{rfc8308}).}
	
\end{itemize}

As a final remark, note that it is the threat model that identifies which subset of the aforementioned attack categories is relevant for a given cryptographic protocol.